% Options for packages loaded elsewhere
\PassOptionsToPackage{unicode}{hyperref}
\PassOptionsToPackage{hyphens}{url}
\PassOptionsToPackage{dvipsnames,svgnames,x11names}{xcolor}
%
\documentclass[
  10pt,
]{article}
\usepackage{amsmath,amssymb}
\usepackage{iftex}
\ifPDFTeX
  \usepackage[T1]{fontenc}
  \usepackage[utf8]{inputenc}
  \usepackage{textcomp} % provide euro and other symbols
\else % if luatex or xetex
  \usepackage{unicode-math} % this also loads fontspec
  \defaultfontfeatures{Scale=MatchLowercase}
  \defaultfontfeatures[\rmfamily]{Ligatures=TeX,Scale=1}
\fi
\usepackage{lmodern}
\ifPDFTeX\else
  % xetex/luatex font selection
\fi
% Use upquote if available, for straight quotes in verbatim environments
\IfFileExists{upquote.sty}{\usepackage{upquote}}{}
\IfFileExists{microtype.sty}{% use microtype if available
  \usepackage[]{microtype}
  \UseMicrotypeSet[protrusion]{basicmath} % disable protrusion for tt fonts
}{}
\makeatletter
\@ifundefined{KOMAClassName}{% if non-KOMA class
  \IfFileExists{parskip.sty}{%
    \usepackage{parskip}
  }{% else
    \setlength{\parindent}{0pt}
    \setlength{\parskip}{6pt plus 2pt minus 1pt}}
}{% if KOMA class
  \KOMAoptions{parskip=half}}
\makeatother
\usepackage{xcolor}
\usepackage[margin=2.2cm]{geometry}
\usepackage{longtable,booktabs,array}
\usepackage{calc} % for calculating minipage widths
% Correct order of tables after \paragraph or \subparagraph
\usepackage{etoolbox}
\makeatletter
\patchcmd\longtable{\par}{\if@noskipsec\mbox{}\fi\par}{}{}
\makeatother
% Allow footnotes in longtable head/foot
\IfFileExists{footnotehyper.sty}{\usepackage{footnotehyper}}{\usepackage{footnote}}
\makesavenoteenv{longtable}
\setlength{\emergencystretch}{3em} % prevent overfull lines
\providecommand{\tightlist}{%
  \setlength{\itemsep}{0pt}\setlength{\parskip}{0pt}}
\setcounter{secnumdepth}{-\maxdimen} % remove section numbering
\sloppy\emergencystretch=3em\usepackage{amssymb}\usepackage{amsmath}\setlength{\tabcolsep}{3pt}\renewcommand{\arraystretch}{1.15}
\ifLuaTeX
  \usepackage{selnolig}  % disable illegal ligatures
\fi
\IfFileExists{bookmark.sty}{\usepackage{bookmark}}{\usepackage{hyperref}}
\IfFileExists{xurl.sty}{\usepackage{xurl}}{} % add URL line breaks if available
\urlstyle{same}
\hypersetup{
  pdftitle={Diameter-7 cells from a spread and its whiskers: the freed sets and the bounds for every q},
  pdfauthor={Degree--diameter lab, 29 October 2026 (job thm-conics)},
  colorlinks=true,
  linkcolor={Maroon},
  filecolor={Maroon},
  citecolor={Blue},
  urlcolor={Blue},
  pdfcreator={LaTeX via pandoc}}

\title{Diameter-7 cells from a spread and its whiskers: the freed sets
and the bounds for every q}
\author{Degree--diameter lab, 29 October 2026 (job thm-conics)}
\date{}

\begin{document}
\maketitle

\emph{Abstract.} The jobs ``conics'' and ``conics-2'' raised the lab's
values at the diameter-7 cells \((14,7)\) to \((20,7)\) by freeing, on
the friendly transitions of the typed hexagon layouts, two sets: the
points of the Hermitian spread \(S\) of \(H(q)\), and the lines meeting
a spread line \(v\) together with an exact whisker set and the \(q^2\)
spread lines that carry no whisker. Here these steps are turned into
statements for every prime power \(q\). (1) For any generalized hexagon
of order \((q,q)\), any spread \(S\), any \(v\in S\) and any whisker set
\(W\), the lines meeting \(v\), the whiskers and the spread lines missed
by the whiskers form a dual segment set; for the exact whisker set of
\texttt{POLISHED.md} (every odd \(q\) and every even \(q\ge4\)) it has
exactly \(q^4+q^2+q\) lines and is maximal, while the spread points form
a segment set of \((q+1)(q^3+1)\) points. (2) For \(2k+1\) levels, any
regular tournament and any typing, some level has at most
\(\lfloor k/2\rfloor\) friendly transitions on each side, which bounds
the hub counts per side; the bound is exact to second order, and the
alternating typing attains it for every \(k\). (3) With the \(q^2+1\)
points of an elliptic quadric section on the non-friendly transitions,
this gives, for every \(k\ge2\) and every prime power \(q\ge q_k\),
\[n_b(q+1+k,7)\ \ge\ 2(2k+1)(q+1)(q^4+q^2+1)+\begin{cases}k\,(q^4+2q^2+q+1)&k\text{ even},\\[2pt]
\tfrac{k-1}2\,(2q^4+q^3+q^2+2q+1)+(k+1)(q^2+1)&k\text{ odd},\end{cases}\]
with \(q_2=16\), \(q_3=q_4=23\), \(q_5=29\), \(q_k\le 300k\) in general.
The hub graphs are given by a new explicit construction (sums of
blown-up hexagons); smaller \(q\) are checked by computer where
possible. The formulas reproduce the vertex counts of all fifteen graphs
built in the two jobs. At \((19,7)\) the formula equals the built value
\(11{,}316{,}940\) and every ingredient is proved.

\emph{Labels.} Every result is marked \textbf{Theorem} (proved here or
in the cited note), \textbf{Computer-checked} (a program checked it; the
program is in the repository) or \textbf{Heuristic}, and
\textbf{{[}new{]}} or \textbf{{[}restatement{]}}. ``New'' means new in
the lab's notes; the literature check of \texttt{D7\_HEXAGON.md} §11
applies, and we did not search the literature again.

\emph{Sources.} \texttt{research/synthetic/CONICS.md} §3.4 (branch
\texttt{claude/conics-generalized-polygons-6gjqxi});
\texttt{research/synthetic/conics2/README.md} (branch
\texttt{claude/conics-2-diameter-7-ic0oat}, state of 11:44 UTC);
\texttt{research/theorems/D7\_HEXAGON.md} (branch
\texttt{claude/d7-hexagon-theorems-4w1prt}), whose Definitions 2.1, 3.1,
Theorems A, B, C, 3.4 and Proposition 5.2 are used throughout;
\texttt{research/theorems/POLISHED.md} §6 (branch
\texttt{claude/thm-polish-2oddll});
\texttt{research/synthetic/SYNTHETIC.md} (branch
\texttt{claude/synthetic-incidence-geometry-s4m08f}), whose conventions
we follow. Programs of this note: \texttt{research/synthetic/d7thm/};
squad round \texttt{research/squads/d7thm-r1/}.

\hypertarget{setting}{%
\section{0. Setting}\label{setting}}

\textbf{Hexagons.} A generalized hexagon \(\Gamma\) of order \((q,q)\)
has \(N=N(q)=(q+1)(q^4+q^2+1)\) points and as many lines; \(q+1\) points
on a line and \(q+1\) lines on a point. Distances are taken in the
incidence graph \(I\) (diameter \(6\), girth \(12\)); elements at
distance \(6\) are \emph{opposite}. For an element \(x\),
\(|\Gamma_i(x)|=(q+1)q^{i-1}\) for \(1\le i\le5\). A \emph{spread} is a
set of \(q^3+1\) pairwise opposite lines; every line not in a spread
\(S\) meets exactly one line of \(S\), in one point
(\texttt{D7\_HEXAGON.md}, Lemma 1.2(c)). The example throughout is the
split Cayley hexagon \(H(q)\) in \(Q(6,q)\) and its Hermitian spread:
the lines of \(H(q)\) inside a hyperplane \(\Pi\) with
\(\Pi\cap Q(6,q)\) elliptic (every \(q\); {[}Th80{]}).

\textbf{Segment sets} (\texttt{D7\_HEXAGON.md}, Definition 1.3). A point
set \(F\) is a \emph{segment set} if for every \(p\in F\) the points of
\(F\setminus\{p\}\) collinear with \(p\) lie on one line through \(p\).
A line set \(G\) is a \emph{dual segment set} if for every \(\ell\in G\)
the lines of \(G\setminus\{\ell\}\) that meet \(\ell\) pass through one
point of \(\ell\). Both kinds have at most \((q+1)(q^3+1)\) elements
(\texttt{D7\_HEXAGON.md}, Proposition 1.4).

\textbf{Layouts.} \(L=2k+1\) \emph{levels}, each a copy of \(I\) of
\emph{type} \(H\) (colour \(\mathcal A\) = points, \(\mathcal B\) =
lines) or \(D\) (colour \(\mathcal A\) = lines, \(\mathcal B\) =
points). The levels are the vertices of a \emph{regular tournament}
\(T\): every two levels are joined by exactly one \emph{transition}
\(j\to j'\), and every level has \(k\) outgoing and \(k\) incoming
transitions. A transition is a bijection from the
\(\mathcal A\)-elements of level \(j\) to the \(\mathcal B\)-elements of
level \(j'\); it is \emph{friendly} if the types differ, and then it is
the identity. \texttt{D7\_HEXAGON.md} uses the rotational tournament
\(j\to j+i\) (\(1\le i\le k\), indices mod \(L\)); Theorem A and Theorem
3.4 there hold verbatim for every regular tournament, since their proofs
only use that every pair of levels is the pair of one transition. The
\emph{freed set} \(R_t\) of a transition is the set of
\(\mathcal A\)-elements whose matching edge is deleted. Level \(j\) then
has \(\alpha_j=\sum_{t\text{ out of }j}|R_t|\) free \(\mathcal A\)-slots
and \(\beta_j=\sum_{t\text{ into }j}|R_t|\) free \(\mathcal B\)-slots.
\(K_{\mathcal B}\le\min_j\alpha_j\) hubs of colour \(\mathcal B\) each
take one free \(\mathcal A\)-slot in every level,
\(K_{\mathcal A}\le\min_j\beta_j\) hubs of colour \(\mathcal A\) one
free \(\mathcal B\)-slot in every level, and a bipartite \emph{hub
graph} \(Y\) joins the two kinds of hubs.

\textbf{What repairs a transition.} Theorem B of
\texttt{D7\_HEXAGON.md}: a friendly transition is repaired if its freed
set is a segment set (source of type \(H\)) or a dual segment set
(source of type \(D\)). Theorem C there, with Corollary 5.3(a): a
non-friendly transition \(H\to H\) with matching \(\sigma\) is repaired
by a set \(R\) of \(m\) pairwise opposite points when \(\sigma\) maps
\(R\) onto \(m\) pairwise opposite lines and the points collinear with
\(x\) onto the lines concurrent with \(\sigma(x)\) for every \(x\in R\);
such a \(\sigma\) exists for every partial ovoid \(R\) and every partial
spread of the same size. A \(D\to D\) transition uses \(\sigma^{-1}\)
and frees \(\sigma(R)\).

\textbf{Hub graphs.} \(\mathrm{HG}(K_1,K_2;\delta)\) denotes a bipartite
graph with colour classes of sizes \(K_1\) and \(K_2\), maximum degree
at most \(\delta\), in which any two vertices \emph{of the same colour}
are at distance at most \(6\).

\textbf{Theorem 0.1} (\texttt{D7\_HEXAGON.md}, Theorem 3.4)
{[}restatement{]}. If every transition is repaired,
\(K_{\mathcal B}\le\min \alpha_j\), \(K_{\mathcal A}\le\min\beta_j\) and
the hub graph is an \(\mathrm{HG}(K_{\mathcal B},K_{\mathcal A};q-k)\),
then the graph has maximum degree \(q+1+k\) and diameter at most \(7\),
and \[n_b(q+1+k,7)\ \ge\ 2(2k+1)N(q)+K_{\mathcal A}+K_{\mathcal B}.\]

The value \(2(2k+1)N(q)\) alone (no freed sets, no hubs) is the value of
the 2013 table at the cells \((7,7)\), \((9,7)\)--\((16,7)\) {[}MS13{]},
and the lab's beyond-table value at \((17,7)\)--\((20,7)\).

\textbf{What uses coordinates.} Sections 2 and 4 are about tournaments
and bipartite graphs only. In §1, Theorem 1.4, Proposition 1.5 and Lemma
1.6 hold in every generalized hexagon of order \((q,q)\) with a spread;
only Theorem 1.7 (the whisker set of \texttt{POLISHED.md} meets each
spread line in \(0\) or \(q+1\) points, so exactly \(q^2\) spread lines
are missed) and the existence statements (Hermitian spread, exact
whisker set, elliptic quadric section) use the coordinates of \(H(q)\).
Coordinates appear only in proofs.

\textbf{Notation for sizes.} \(P(q)=(q+1)(q^3+1)=q^4+q^3+q+1\),
\(\Lambda(q)=q^4+q^2+q\), and \(m\) the size of the non-friendly sets.
Note \(P-\Lambda=q^3-q^2+1\).

\hypertarget{the-freed-sets}{%
\section{1. The freed sets}\label{the-freed-sets}}

\hypertarget{the-configuration}{%
\subsection{1.1 The configuration}\label{the-configuration}}

\textbf{Definition 1.1.} Let \(S\) be a spread of \(\Gamma\),
\(\mathcal P_S\) the set of points on its lines, and \(v\in S\).
\(\mathcal L_v\) is the set of the \(q(q+1)\) lines other than \(v\)
that meet \(v\). A point is \emph{free} if it is at distance \(5\) from
\(v\) and not in \(\mathcal P_S\). A \emph{whisker set} is a set \(W\)
of lines, each opposite \(v\) and not in \(S\), such that no free point
lies on two lines of \(W\); it is \emph{exact} if every free point lies
on exactly one line of \(W\). Finally
\[S_0(W)=\{n\in S\setminus\{v\}:\ n\text{ meets no line of }W\},\qquad R_L(W)=\mathcal L_v\cup W\cup S_0(W).\]

(This is the configuration of Theorem D, \texttt{POLISHED.md} §6.2 and
\texttt{SYNTHETIC.md} §3.2, written without coordinates.)

\textbf{Lemma 1.2} {[}restatement, with the counts; Theorem{]}. (a) The
points at distance at most \(3\) from \(v\) are the \(q+1\) points of
\(v\) and the \(q^2(q+1)\) points at distance \(3\), each of which lies
on exactly one line of \(\mathcal L_v\). (b) The points of
\(\mathcal P_S\) not on \(v\) are at distance \(5\) from \(v\). (c)
There are \(q^3(q^2-1)\) free points. (d) A line opposite \(v\) and not
in \(S\) has exactly one point in \(\mathcal P_S\) and \(q\) free
points. (e) A whisker set has at most \(q^2(q^2-1)\) lines, with
equality if and only if it is exact.

\emph{Proof.} (a) A point at distance \(3\) from \(v\) lies on a line
meeting \(v\); two such lines through it would close a circuit of length
\(4\) or \(6\). Counting: \(|\Gamma_3(v)|=(q+1)q^2\). (b) A point \(x\)
on a spread line \(n\ne v\) has \(d(x,v)\) odd, at most \(5\) and at
least \(d(n,v)-1=5\). (c) \(|\Gamma_5(v)|=(q+1)q^4\), of which
\(q^3(q+1)\) lie in \(\mathcal P_S\) by (b). (d) Every point of a line
opposite \(v\) is at distance \(5\) from \(v\); the line meets exactly
one line of \(S\), in one point. (e) By (d) and (c), counting incidences
between free points and \(W\). \(\square\)

\hypertarget{spread-points-and-whiskers}{%
\subsection{1.2 Spread points and
whiskers}\label{spread-points-and-whiskers}}

\textbf{Proposition 1.3} (\texttt{D7\_HEXAGON.md}, Proposition 1.4)
{[}restatement; Theorem{]}. \(\mathcal P_S\) is a segment set of
\(P(q)=(q+1)(q^3+1)\) points, and no segment set is larger; equality
holds only for the point sets of spreads.

\textbf{Theorem 1.4} {[}new in this form; Theorem{]}. Let \(S\) be a
spread, \(v\in S\) and \(W\) any whisker set.

\begin{enumerate}
\def\labelenumi{(\alph{enumi})}
\item
  \(R_L(W)=\mathcal L_v\cup W\cup S_0(W)\) is a dual segment set, and
  two of its lines that meet, meet in a point of \(\mathcal P_S\).
\item
  \(|R_L(W)|=q(q+1)+|W|+|S_0(W)|\).
\item
  If \(W\) is exact and every point of \(\mathcal P_S\) not on \(v\)
  lies on a line of \(W\) or of \(S_0(W)\), then every point of
  \(\Gamma\) lies on a line of \(R_L(W)\), and no line can be added to
  \(R_L(W)\): it is a maximal dual segment set.
\item
  A line \(n\in S\setminus\{v\}\) all of whose whiskers pass through one
  point of \(n\) can be added to \(R_L(W)\) as well.
\end{enumerate}

\emph{Proof.} (a) A line of \(\mathcal L_v\) has one point on
\(v\subseteq\mathcal P_S\) and its other \(q\) points at distance \(3\)
from \(v\), so not in \(\mathcal P_S\) (Lemma 1.2(b)). A whisker has one
point in \(\mathcal P_S\) and \(q\) free points (Lemma 1.2(d)). So every
line of \(\mathcal L_v\cup W\) has exactly one point in
\(\mathcal P_S\). Its other points are either at distance \(3\) from
\(v\), and then on no other line of \(\mathcal L_v\) (Lemma 1.2(a)) and
on no whisker (whose points are at distance \(5\)), or free, and then on
no line of \(\mathcal L_v\) and on at most one whisker. So two lines of
\(\mathcal L_v\cup W\) never meet outside \(\mathcal P_S\). A line
\(n\in S_0(W)\) meets no line of \(\mathcal L_v\) (whose points are at
distance at most \(3\) from \(v\), while those of \(n\) are at distance
\(5\)), no whisker (definition) and no other spread line. So the lines
of \(S_0(W)\) are isolated, and for every other line \(\ell\) of
\(R_L(W)\) the lines of \(R_L(W)\) meeting \(\ell\) all pass through the
unique point of \(\ell\) in \(\mathcal P_S\). (b) The three sets are
disjoint: whiskers are opposite \(v\), lines of \(\mathcal L_v\) are
not, and neither kind is in \(S\). (c) Points of \(v\) and points at
distance \(3\) from \(v\) lie on lines of \(\mathcal L_v\), free points
on whiskers (\(W\) exact), and the points of \(\mathcal P_S\) off \(v\)
by hypothesis. A line \(\ell\notin R_L(W)\) then meets a line of
\(R_L(W)\) at each of its \(q+1\) points, and these lines are different
from each other (two points of \(\ell\) on one line \(\ne\ell\) is
impossible), so \(R_L(W)\cup\{\ell\}\) is not a dual segment set. (d)
Such an \(n\) meets lines of \(R_L(W)\) only in that point \(x\) (it
meets no line of \(\mathcal L_v\) and no other spread line), and for a
whisker through \(x\) the new line meets it at its point in
\(\mathcal P_S\), where all its other meetings already are. Two such
spread lines do not meet. \(\square\)

For \(W\) exact and \(S_0(W)\) left out this is \texttt{CONICS.md},
Lemma 3.1; for the whisker set of \texttt{POLISHED.md} and its \(q^2\)
empty spread lines it is Proposition 1(c),(d) of
\texttt{conics2/README.md}. The proof above uses only incidences.

\textbf{Proposition 1.5} {[}new; Theorem{]} (\emph{where the room is}).
Let \(G\) be a dual segment set in which any two lines that meet do so
in a point of \(\mathcal P_S\) (as in Theorem 1.4(a)). Then
\(|G\setminus S|\le q^4+q\), with equality if and only if the lines of
\(G\setminus S\) cover every point outside \(\mathcal P_S\) exactly
once. Hence \(|G|\le q^4+q+|G\cap S|\). For an exact whisker set,
\(\mathcal L_v\cup W\) has exactly \(q^4+q\) lines.

\emph{Proof.} A line not in \(S\) meets exactly one spread line, in one
point, so it has exactly one point in \(\mathcal P_S\) (a second point
would lie on a second spread line). Its other \(q\) points lie outside
\(\mathcal P_S\), and by hypothesis these \(q\)-sets are disjoint for
different lines of \(G\setminus S\). There are
\(N-P(q)=q^5+q^2=q(q^4+q)\) points outside \(\mathcal P_S\). Finally
\(q(q+1)+q^2(q^2-1)=q^4+q\). \(\square\)

So within this class \(R_L(W)\) can only grow by containing more spread
lines, at most \(q^3-q^2+1\) more; the general bound
\((q+1)(q^3+1)=q^4+q+(q^3+1)\) is the case where all spread lines are in
\(G\), which forces an ovoid.

\textbf{Lemma 1.6} {[}new; Theorem{]}. Let \(W\) be exact. Every whisker
meets exactly one line of \(S\setminus\{v\}\), and a line
\(n\in S\setminus\{v\}\) meets at most \(q^2-1\) whiskers. So at least
\(q^2\) spread lines other than \(v\) meet a whisker and
\(|S_0(W)|\le q^3-q^2\). If every line of \(S\setminus\{v\}\) meets
either no whisker or \(q+1\) whiskers, then \(|S_0(W)|=q^2\).

\emph{Proof.} A whisker meets one line of \(S\) (§0), not \(v\), since
it is opposite \(v\). At a point \(x\) of \(n\) the \(q+1\) lines are
\(n\), the line through \(x\) on the geodesic to \(v\) (at distance
\(4\) from \(v\), so not a whisker) and \(q-1\) others. Hence \(n\)
meets at most \((q+1)(q-1)\) whiskers, and \(q^2(q^2-1)/(q^2-1)=q^2\).
In the last case the lines meeting whiskers number
\(q^2(q^2-1)/(q+1)=q^2(q-1)\), and \(q^3-q^2(q-1)=q^2\). \(\square\)

\hypertarget{the-exact-whisker-set-of-polished.md}{%
\subsection{\texorpdfstring{1.3 The exact whisker set of
\texttt{POLISHED.md}}{1.3 The exact whisker set of POLISHED.md}}\label{the-exact-whisker-set-of-polished.md}}

In \(H(q)\) with the coordinates of \texttt{POLISHED.md} §6.1 (\(q\)
odd) and §6.5 (\(q\) even): \(\Pi\) the elliptic hyperplane
\(X_0+cX_4=0\) (\(q\) odd, \(-c\) a non-square) or
\(X_3=X_0+\delta X_4\) (\(q\) even, \(\mathrm{Tr}\,\delta=1\)), \(S\)
its Hermitian spread, \(v=\langle e_1,e_6\rangle\), the hyperplanes
\(\Pi_z\supset v\) (\(z=(\alpha,\beta)\in\mathbb F^2\)) and their
spreads \(S_z\), the function \(\nu\) on the points of \(\Pi\) off
\(v\), the zero-free cubic \(G\) and \(F(z)=G(z)/N(z)\) with \(N\) the
anisotropic norm form;
\(W=\{\ell\in S_z\setminus\{v\}: z\neq0,\ \nu(\ell\cap\Pi)=F(z)\}\) is
exact for every odd \(q\) (Theorem 6.5) and every even \(q\ge4\)
(Theorem 6.14).

\textbf{Theorem 1.7} (\texttt{conics2/README.md}, Proposition 1(a),(b))
{[}restatement with proof; Theorem{]}. Let \(q\) be odd, or even with
\(q\ge4\), and \(W\) the exact whisker set above. Then every point \(X\)
of \(\mathcal P_S\) not on \(v\) lies on exactly one whisker if
\(\nu(X)\neq0\) and on none if \(\nu(X)=0\); \(\nu\) is constant on
every line of \(S\setminus\{v\}\). Hence every line of
\(S\setminus\{v\}\) meets either \(q+1\) whiskers, one at each of its
points, or none, and \(S_0(W)=\{n\in S\setminus\{v\}:\nu(n)=0\}\) has
exactly \(q^2\) lines.

\emph{Proof.} Whiskers through \(X\) are whiskers whose point in \(\Pi\)
is \(X\) (their other points are off \(\Pi\)). On \(\Pi\) the equation
of \(\Pi_z\) reduces to \(\alpha X_2+\beta X_5=0\) (both parities), and
\((X_2,X_5)\neq(0,0)\) because \(X\) is not on \(v\)
(\texttt{POLISHED.md}, Lemma 6.1(d) / 6.10). So the \(z\neq0\) with
\(X\in\Pi_z\) form the punctured line \(\{tz_0:t\neq0\}\),
\(z_0=(X_5,-X_2)\). For each of them exactly one line of \(S_z\) passes
through \(X\) (Lemma 6.1(b); for even \(q\) read Lemma 6.10 for 6.1
throughout); it is not \(v\), it is not in \(S\) (6.1(e)), so it meets
\(\Pi\) only in \(X\) (6.1(c)), and different \(z\) give different lines
(6.1(e)). It is a whisker if and only if \(\nu(X)=F(tz_0)\). As \(G\) is
a cubic form and \(N\) a quadratic form, \(F(tz_0)=tF(z_0)\), and
\(F(z_0)\neq0\) since \(G\) has no zero. So \(t\mapsto F(tz_0)\) is a
bijection of \(\mathbb F^\times\): one whisker through \(X\) if
\(\nu(X)\neq0\), none if \(\nu(X)=0\). \(\nu\) is constant on spread
lines by \texttt{POLISHED.md}, Lemma 6.6(d) / 6.12(d). The count is
Lemma 1.6. \(\square\)

\textbf{Corollary 1.8} {[}Theorem{]}. For every prime power \(q\) with
\(q\) odd or \(q\ge4\), in \(H(q)\): \(R_H=\mathcal P_S\) is a segment
set with \(P(q)=q^4+q^3+q+1\) points, and
\(R_L=\mathcal L_v\cup W\cup S_0(W)\) is a maximal dual segment set with
\(\Lambda(q)=q^4+q^2+q\) lines (Theorem 1.4 with Theorem 1.7). In the
notation of Proposition 1.5, \(\mathcal L_v\cup W\) is extremal and
\(|R_L\cap S|=q^2\).

\emph{Computer-checked.} \texttt{conics2/typed\_build.py} asserts every
statement of Theorem 1.7 and Corollary 1.8 for \(q=4,5,7,11,13,16\)
inside the builds. Independently, squad round \texttt{d7thm-r1} (own
code, not the builder): card c1 for \(q=3,5,7\) (hexagon built from the
quadric and Tits' equations), card c2 for \(q=9,11\), card c3 for
\(q=2,4,8\): exactness of \(W\), \(\nu\) constant on spread lines, \(0\)
or \(q+1\) whiskers per spread line, \(|S_0|=q^2\) equal to the lines
with \(\nu=0\), the dual segment property, \(|R_L|=q^4+q^2+q\)
(\(6{,}651\) at \(q=9\), \(14{,}773\) at \(q=11\), \(276\) at \(q=4\),
\(4{,}168\) at \(q=8\)), maximality (every outside line tested), and the
segment property of \(\mathcal P_S\); card c2 also for every admissible
\(c\) at \(q=9,11\), card c3 for every admissible \((c_0,c_1)\) at
\(q=4,8\); negative controls fail as they should
(\texttt{research/squads/d7thm-r1/scoreboard.jsonl}).

\emph{Remarks.} (i) \(q=2\): Lemma 6.11 of \texttt{POLISHED.md} gives no
zero-free \(G\) (card c3 re-checked all four pairs), but an exact
whisker set of \(12\) lines exists (\texttt{POLISHED.md}, Corollary
6.15), so Theorem 1.4 gives \(|R_L|\ge18\); the largest dual segment set
of \(H(2)\) has \(25\) lines (\texttt{conics2/README.md} §3, SAT). (ii)
\(R_L\) is not maximum in general: \(H(4)\) has a dual segment set of
\(288>276\) lines (same place; it is not of the kind in Proposition
1.5). Whether \(q^4+q^2+q\) can be beaten for \(q\ge5\) is open; by
Proposition 1.5 a set of this kind would need more than \(q^2\) spread
lines. (iii) For \(q=3^h\) the pencils of an ovoid give a dual segment
set of \(P(q)\) lines, and the duality layout of
\texttt{D7\_HEXAGON.md}, Corollary 8.2 is better than any typed layout.
(iv) We know no coordinate-free description of the \(q^2\) lines with
\(\nu=0\).

\hypertarget{the-non-friendly-sets}{%
\subsection{1.4 The non-friendly sets}\label{the-non-friendly-sets}}

\textbf{Proposition 1.9} (\texttt{D7\_HEXAGON.md}, Lemma 1.2(a),
Corollaries 5.3(a) and 9.1) {[}restatement; Theorem{]}. For every prime
power \(q\), the \(q^2+1\) points of an elliptic quadric section
\(Q^-(3,q)\) of \(Q(6,q)\) are pairwise opposite in \(H(q)\), and with
any \(q^2+1\) lines of the Hermitian spread they give a local duality
\(\sigma\) as in Theorem C. So \(m=q^2+1\) is available on every
non-friendly transition.

The lab's builds use larger non-friendly sets found by computer at
\(q=11\) (\(m=768\), mirrored set) and \(q=13\) (\(m=792\)), and
\(m=q^2+1=257\) at \(q=16\) (\textbf{Computer-checked}: the builder
re-checks non-collinearity, non-concurrency and the local duality at
every point).

\hypertarget{slot-counts}{%
\section{2. Slot counts}\label{slot-counts}}

Let \(T\) be a regular tournament on \(2k+1\) levels with a typing. For
a level \(j\) let \(f_j\) be the number of friendly transitions out of
\(j\) and \(g_j\) the number into \(j\). With friendly sets of sizes
\(P\) (points, source \(H\)) and \(\Lambda\) (lines, source \(D\)) and
non-friendly sets of size \(m\),
\[\alpha_j=f_jX_j+(k-f_j)m,\quad X_j=\begin{cases}P&\tau_j=H\\ \Lambda&\tau_j=D\end{cases},\qquad
\beta_j=g_jY_j+(k-g_j)m,\quad Y_j=\begin{cases}\Lambda&\tau_j=H\\ P&\tau_j=D\end{cases}\]
(a friendly transition into an \(H\) level comes from a \(D\) level,
which frees lines). Put \(a=\lfloor k/2\rfloor\) and assume
\[P\ \ge\ \Lambda\ \ge\ m,\qquad \Lambda-m\ \ge\ a\,(P-\Lambda).\tag{$\ast$}\]
For the sets of §1, \((\ast)\) holds whenever \(k\le 2q+1\) and
\(m\le q^2+1\), since then \(a(q^3-q^2+1)\le q^4-q^3+q\le q^4+q-1\); at
\(q=11\), \(m=768\) it holds for \(k\le23\), at \(q=13\), \(m=792\) for
\(k\le27\). The bound needs \(k\le q-1\) anyway (hub degree
\(q-k\ge1\)).

\textbf{Theorem 2.1} {[}new; Theorem{]}. For every regular tournament on
\(2k+1\) levels and every typing, with \(a=\lfloor k/2\rfloor\):

\begin{enumerate}
\def\labelenumi{(\alph{enumi})}
\item
  \(\min_jf_j\le a\) and \(\min_jg_j\le a\);
\item
  if \(k\) is even, it is impossible that \(f_j\ge a\) at every level
  and \(f_j\ge a+1\) at every \(D\) level; likewise with \(g\) and \(H\)
  levels;
\item
  if \(k\) is odd, it is impossible that \(f_j\ge a+1\) at every \(D\)
  level and \(g_j\ge a+1\) at every \(H\) level.
\end{enumerate}

Consequently, if \(P\ge\Lambda\ge m\): for \(k\) odd,
\(\min\alpha,\min\beta\le aP+(a+1)m\) and
\(\min\alpha+\min\beta\le a(P+\Lambda)+2(a+1)m\); for \(k\) even and
under \((\ast)\), \(\min\alpha\le a(\Lambda+m)\) and
\(\min\beta\le a(\Lambda+m)\).

\emph{Proof.} Let \(h\) and \(d\) be the numbers of \(H\) and \(D\)
levels. The \(k-f_j\) non-friendly transitions out of an \(H\) level go
to \(H\) levels, so \(\sum_{\tau_j=H}(k-f_j)=\binom h2\), and likewise
\(\sum_{\tau_j=D}(k-f_j)=\binom d2\) and the same with \(g\).

\begin{enumerate}
\def\labelenumi{(\alph{enumi})}
\item
  If \(f_j\ge a+1\) everywhere, then \(\binom h2\le h(k-a-1)\) and
  \(\binom d2\le d(k-a-1)\), so \(h+d\le 4k-4a-2\le 2k\), less than
  \(2k+1\). For \(g\) the same.
\item
  \(k=2a\). From \(f_j\ge a\) on \(H\) levels, \(\binom h2\le ha\), so
  \(h\le 2a+1\); from \(f_j\ge a+1\) on \(D\) levels,
  \(\binom d2\le d(a-1)\), so \(d\le 2a-1\). Then \(h+d\le4a<4a+1\). For
  \(g\) exchange the roles of \(H\) and \(D\).
\item
  \(k=2a+1\). On \(H\) levels \(g_j\ge a+1\) gives \(\binom h2\le ha\)
  (count non-friendly incoming transitions), so \(h\le2a+1\); on \(D\)
  levels \(f_j\ge a+1\) gives \(\binom d2\le da\), so \(d\le 2a+1\).
  Then \(h+d\le4a+2<4a+3\).
\end{enumerate}

Consequences. A level with \(f_j\le a-1\) has
\(\alpha_j\le(a-1)P+(a+1)m\le a\Lambda+am\), by \((\ast)\). If \(k\) is
even and \(\min\alpha>a(\Lambda+m)\), every level has \(f_j\ge a\), and
every \(D\) level \(f_j\ge a+1\) (a \(D\) level with \(f_j=a\) has
\(\alpha_j=a\Lambda+am\)), contradicting (b). If \(k\) is odd, a level
with \(f_j\le a\) has \(\alpha_j\le aP+(a+1)m\), so by (a)
\(\min\alpha\le aP+(a+1)m\), and the same for \(\beta\). If
\(\min\alpha+\min\beta>a(P+\Lambda)+2(a+1)m\), both minima exceed
\(a\Lambda+(a+1)m\). A \(D\) level with \(f_j\le a\) has
\(\alpha_j\le a\Lambda+(a+1)m\) and an \(H\) level with \(g_j\le a\) has
\(\beta_j\le a\Lambda+(a+1)m\), so every \(D\) level has \(f_j\ge a+1\)
and every \(H\) level \(g_j\ge a+1\), contradicting (c). \(\square\)

\textbf{Theorem 2.2} {[}new; Theorem{]} (\emph{the alternating layout}).
Take the rotational tournament \(j\to j+i\) (\(1\le i\le k\)) on
\(\mathbb Z_{2k+1}\) and the typing \(\tau_j=H\) for \(j\) even,
\(\tau_j=D\) for \(j\) odd (\(0\le j\le 2k\); the levels \(2k\) and
\(0\) are both of type \(H\)). Then \(f_j=\lceil k/2\rceil\) for
\(0\le j\le k\), and for \(j=k+t\) (\(1\le t\le k\))
\[f_j=\begin{cases}k/2\ (t\text{ even, }H)\quad k/2+1\ (t\text{ odd, }D)&k\text{ even},\\
(k-1)/2\ (t\text{ odd, }H)\quad (k+1)/2\ (t\text{ even, }D)&k\text{ odd};\end{cases}\]
and \(g_j=f_{2k-j}\). Hence, under \((\ast)\),
\[k\text{ even: }\ \min\alpha=\min\beta=\tfrac k2(\Lambda+m);\qquad k\text{ odd: }\ \min\alpha=\tfrac{k-1}2P+\tfrac{k+1}2m,
\quad\min\beta=\tfrac{k-1}2\Lambda+\tfrac{k+1}2m,\] which are the bounds
of Theorem 2.1.

\emph{Proof.} For \(j\le k\) the targets \(j+1,\dots,j+k\) need no
reduction mod \(2k+1\), and exactly \(\lceil k/2\rceil\) of them have
the parity opposite to \(j\). For \(j=k+t\) the targets are
\(j+1,\dots,2k\) (\(k-t\) of them, \(\lceil (k-t)/2\rceil\) of the other
parity) and \(0,\dots,t-1\), of which \(\lfloor t/2\rfloor\) (if \(j\)
is even) or \(\lceil t/2\rceil\) (if \(j\) is odd) have the other
parity; the four cases give the values stated. The reflection
\(j\mapsto 2k-j\) preserves types and reverses every transition, so
\(g_j=f_{2k-j}\). For \(k\) even the minimum \(a\) of \(f\) is attained
at \(D\) levels (odd \(j\le k\)), so \(\min\alpha=a\Lambda+am\); the
minimum of \(g\) is attained at \(H\) levels, where friendly incoming
sets are lines: \(\min\beta=a\Lambda+am\). For \(k\) odd, \(f_j=a\) only
at \(H\) levels and every \(D\) level has \(f_j=a+1\), whose
\(\alpha_j=(a+1)\Lambda+am\ge aP+(a+1)m\) by \((\ast)\); so
\(\min\alpha=aP+(a+1)m\); \(g_j=a\) only at \(H\) levels, where
\(\beta_j=a\Lambda+(a+1)m\). \(\square\)

So for the hubs the alternating layout is optimal among all regular
tournaments and typings, and the slot counts are
\[K_{\mathcal B}(q,k)=K_{\mathcal A}(q,k)=\tfrac k2(\Lambda+m)\ (k\text{ even}),\qquad
K_{\mathcal B}=\tfrac{k-1}2P+\tfrac{k+1}2m,\ \ K_{\mathcal A}=\tfrac{k-1}2\Lambda+\tfrac{k+1}2m\ (k\text{ odd}).\]
For \(k=1\) this is \(m\) on both sides, \texttt{D7\_HEXAGON.md}
Corollary 9.1.

\emph{The layouts of the builds} {[}Computer-checked against the
builders' logs{]}. The free slots per level, in level order:

\begin{longtable}[]{@{}
  >{\raggedright\arraybackslash}p{(\columnwidth - 6\tabcolsep) * \real{0.2500}}
  >{\raggedright\arraybackslash}p{(\columnwidth - 6\tabcolsep) * \real{0.2500}}
  >{\raggedright\arraybackslash}p{(\columnwidth - 6\tabcolsep) * \real{0.2500}}
  >{\raggedright\arraybackslash}p{(\columnwidth - 6\tabcolsep) * \real{0.2500}}@{}}
\toprule\noalign{}
\begin{minipage}[b]{\linewidth}\raggedright
\(k\)
\end{minipage} & \begin{minipage}[b]{\linewidth}\raggedright
types
\end{minipage} & \begin{minipage}[b]{\linewidth}\raggedright
\(\alpha_j\)
\end{minipage} & \begin{minipage}[b]{\linewidth}\raggedright
\(\beta_j\)
\end{minipage} \\
\midrule\noalign{}
\endhead
\bottomrule\noalign{}
\endlastfoot
2 & \(HDHDH\) (alternating) & \(P{+}m\), \(\Lambda{+}m\), \(P{+}m\),
\(2\Lambda\), \(P{+}m\) & \(\Lambda{+}m\), \(2P\), \(\Lambda{+}m\),
\(P{+}m\), \(\Lambda{+}m\) \\
3 & \(HHDHHDD\) & \(P{+}2m\), \(P{+}2m\), \(2\Lambda{+}m\), \(2P{+}m\),
\(2P{+}m\), \(2\Lambda{+}m\), \(2\Lambda{+}m\) & \(2\Lambda{+}m\),
\(2\Lambda{+}m\), \(2P{+}m\), \(\Lambda{+}2m\), \(\Lambda{+}2m\),
\(2P{+}m\), \(2P{+}m\) \\
4 & \(HHDHDHDHD\) (alternating, rotated) & \(2P{+}2m\), \(2P{+}2m\),
\(2\Lambda{+}2m\), \(2P{+}2m\), \(2\Lambda{+}2m\), \(2P{+}2m\),
\(3\Lambda{+}m\), \(2P{+}2m\), \(3\Lambda{+}m\) & \(2\Lambda{+}2m\),
\(2\Lambda{+}2m\), \(3P{+}m\), \(2\Lambda{+}2m\), \(3P{+}m\),
\(2\Lambda{+}2m\), \(2P{+}2m\), \(2\Lambda{+}2m\), \(2P{+}2m\) \\
\end{longtable}

The \(k=3\) layout of the builds is not the alternating one but has the
same minima. With the values of each build these vectors equal, entry by
entry, the vectors printed by \texttt{typed\_build.py}
(\texttt{d7thm/cells.out}).

\emph{Computer-checked.} \texttt{d7thm/slots.py}: for \(k\le8\) every
typing of the rotational tournament (the optimum is the one of Theorem
2.2, numerically for \(q=11\) with \(m=768\), \(q=13\) with \(m=792\),
\(q=16,101,1009\)); every labelled regular tournament on \(5\) and \(7\)
levels (\(24\) and \(2{,}640\)) with every typing; SAT (CaDiCaL) on
\(9,11,13\) levels confirms the impossibility statements. Independently
squad card c6: all regular tournaments on \(3,5,7\) levels,
\(200{,}000\) random ones on \(9\) and \(11\) levels, annealing, and the
closed forms of Theorem 2.2 for \(k\le40\). Card c6 also found that
without \((\ast)\) the value of \(\min\alpha\) for odd \(k\) changes
(the \(D\) levels become the minimum); this happens only for \(k\) far
above \(q\) (for instance \(k\ge27\) at \(q=11\)), outside the range of
the bound.

\hypertarget{the-bound}{%
\section{3. The bound}\label{the-bound}}

\textbf{Theorem 3.1} {[}new; Theorem{]}. Let \(q\) be a prime power,
\(q\) odd or \(q\ge4\), let \(1\le k\le q-1\), and let
\(K_{\mathcal B}(q,k)\), \(K_{\mathcal A}(q,k)\) be the slot counts of
Theorem 2.2 with \(P=P(q)\), \(\Lambda=\Lambda(q)\) and \(m=q^2+1\). If
an \(\mathrm{HG}(K_1,K_2;q-k)\) exists with \(K_1\le K_{\mathcal B}\)
and \(K_2\le K_{\mathcal A}\), then
\[n_b(q+1+k,\,7)\ \ge\ 2(2k+1)N(q)+K_1+K_2 .\]

\emph{Proof.} Take \(2k+1\) copies of the incidence graph of \(H(q)\),
the alternating layout of Theorem 2.2, the identity on friendly
transitions with freed sets \(R_H=\mathcal P_S\) (source \(H\)) and
\(R_L\) (source \(D\)) of Corollary 1.8, and on non-friendly transitions
the local duality \(\sigma\) (\(H\to H\)) or \(\sigma^{-1}\)
(\(D\to D\)) of Proposition 1.9 with the \(q^2+1\) freed elements. Every
friendly transition is repaired by Theorem B, every other one by Theorem
C. Attach \(K_1\) hubs of colour \(\mathcal B\) and \(K_2\) of colour
\(\mathcal A\) to free slots in any way (one slot per level each; unused
slots stay free) and use the given graph as hub graph. Theorem 0.1
applies. \(\square\)

The same holds with \(m\) replaced by the size of any non-friendly set
satisfying Theorem C, provided \((\ast)\) holds; and for any generalized
hexagon of order \((q,q)\) with a spread, an exact whisker set with the
property of Theorem 1.7, and a partial ovoid, by the same proof.

\textbf{Corollary 3.2} {[}new; Theorem when the hub graph exists{]}.
With full slots (\(K_1=K_{\mathcal B}\), \(K_2=K_{\mathcal A}\)) Theorem
3.1 gives \(n_b(q+1+k,7)\ge F_k(q)\), where
\[F_k(q)=2(2k+1)(q+1)(q^4+q^2+1)+\begin{cases}k\,(q^4+2q^2+q+1),&k\text{ even},\\[2pt]
\tfrac{k-1}2\,(2q^4+q^3+q^2+2q+1)+(k+1)(q^2+1),&k\text{ odd}.\end{cases}\]
In particular
\[F_2(q)=10N(q)+2q^4+4q^2+2q+2,\quad F_3(q)=14N(q)+2q^4+q^3+5q^2+2q+5,\quad F_4(q)=18N(q)+4q^4+8q^2+4q+4.\]
The gain over the 2013 table value \(2(2k+1)N(q)\) is \(kq^4+O(q^2)\)
for \(k\) even and \((k-1)q^4+\frac{k-1}2q^3+O(q^2)\) for \(k\) odd; the
earlier general statement, \texttt{D7\_HEXAGON.md} Corollary 9.3
(\(k=2\)), gave \(2(q+2)(q^2+1)\approx2q^3\), and none was stated there
for \(k\ge3\). With a larger non-friendly set \(m\) the formula changes
by \(k(m-q^2-1)\) for \(k\) even and \((k+1)(m-q^2-1)\) for \(k\) odd.

\emph{Computer-checked} (\texttt{d7thm/cells.py}). The closed forms
agree with the slot formulas for \(q\le37\), \(k\le6\).

\hypertarget{hub-graphs}{%
\section{4. Hub graphs}\label{hub-graphs}}

The bound needs \(\mathrm{HG}(K_{\mathcal B},K_{\mathcal A};q-k)\) with
\(K\approx\lfloor k/2\rfloor q^4\) hubs per side and degree \(q-k\). The
Moore-type bound for same-colour distance \(6\) is about \((q-k)^6\), so
there is room, but the explicit hub graph of \texttt{D7\_HEXAGON.md}
(Lemma 7.1, a plane with pendant trees) has only about
\(\frac4{27}(q-k)^4\) vertices per colour. Write
\(K_0(s)=(s+1)(s^4+s^2+1)\) for the number of points of a hexagon of
order \((s,s)\).

\textbf{Lemma 4.1} (\emph{blown-up hexagon}) {[}Theorem; elementary{]}.
Let \(s\) be a prime power and \(c\ge1\). Replace every point and every
line of \(H(s)\) by between \(1\) and \(c\) copies, and join two copies
when the originals are incident. The result is an
\(\mathrm{HG}(K_1,K_2;c(s+1))\) for any chosen
\(K_1,K_2\in[K_0(s),cK_0(s)]\), and any two of its vertices of different
colours are at distance at most \(5\).

\emph{Proof.} A path between copies of \(x\neq y\) projects to a walk
from \(x\) to \(y\) of the same length, and a geodesic from \(x\) to
\(y\) lifts to copies, so the distance is \(d(x,y)\); two copies of one
element are at distance \(2\). A copy of \(x\) has degree
\(\sum_{y\ \mathrm I\ x}t_y\le c(s+1)\). \(\square\)

\textbf{Lemma 4.2} (\emph{sums}) {[}new; Theorem{]}. Let
\(Y_1,\dots,Y_r\) be bipartite graphs with colour classes \((A_i,B_i)\),
each with all same-colour distances at most \(6\), and for \(i<r\) also
all different-colour distances at most \(5\). Suppose \(|A_j|\le|B_i|\)
and \(|B_j|\le|A_i|\) for \(i<j\). For every \(i<j\) add an injective
map \(A_j\to B_i\) and an injective map \(B_j\to A_i\), as edges. The
result is an
\(\mathrm{HG}(\sum|A_i|,\sum|B_i|;\ \max_i\Delta(Y_i)+r-1)\).

\emph{Proof.} The new edges join the colours \(A\) and \(B\), so the
graph is bipartite with classes \(\bigcup A_i\), \(\bigcup B_i\). A
vertex of \(Y_i\) gains one edge to each \(Y_j\), \(j<i\), and at most
one from each \(Y_j\), \(j>i\). Same-colour pairs inside one \(Y_i\) are
at distance at most \(6\) in \(Y_i\). For \(x\in A_j\) and \(z\in A_i\)
with \(i<j\), \(x\) has a neighbour \(b\in B_i\), and
\(d_{Y_i}(b,z)\le5\) since \(i<r\); so \(d(x,z)\le6\). The same for
\(B\). \(\square\)

\textbf{Lemma 4.3} (\texttt{D7\_HEXAGON.md}, Lemma 7.1, with the colours
sized separately) {[}restatement; Theorem{]}. Let \(s\) be a prime
power, \(n=s^2+s+1\), \(t=\delta-s-1\ge0\), \(u=\delta-1\). Attach to
the incidence graph of \(\mathrm{PG}(2,s)\) \(c_P\) pendant vertices at
the points and \(c_L\) at the lines (at most \(t\) per core vertex), and
\(g_P\), \(g_L\) pendant vertices below these (at most \(u\) per
vertex). This is an \(\mathrm{HG}(n+c_L+g_P,\ n+c_P+g_L;\ \delta)\).
\emph{Proof.} The distance argument of Lemma 7.1 does not use the
numbers of pendant vertices. \(\square\) (Different-colour distances can
be \(7\) here, so such a part can only be the last part \(Y_r\) in Lemma
4.2; squad card c5 confirmed this with a negative control.)

\textbf{Theorem 4.4} {[}new; Theorem{]}. Let \(k\ge2\). An
\(\mathrm{HG}(K_{\mathcal B}(q,k),K_{\mathcal A}(q,k);q-k)\) (full slots
of Theorem 2.2, \(m=q^2+1\)) exists

\begin{enumerate}
\def\labelenumi{(\alph{enumi})}
\item
  for every prime power \(q\ge300k\);
\item
  for \(2\le k\le12\), for every prime power \(q\ge q_k\), where
\end{enumerate}

\begin{longtable}[]{@{}llllllllllll@{}}
\toprule\noalign{}
\(k\) & 2 & 3 & 4 & 5 & 6 & 7 & 8 & 9 & 10 & 11 & 12 \\
\midrule\noalign{}
\endhead
\bottomrule\noalign{}
\endlastfoot
\(q_k\) & 16 & 23 & 23 & 29 & 23 & 37 & 29 & 41 & 32 & 49 & 41 \\
\end{longtable}

(For \(k=10\) every prime power \(q\ge 32\); the values are not monotone
in \(k\) because prime powers are irregular.)

\emph{Proof.} (a) Let \(q\ge300k\) (so \(q\ge600\)). The slot counts
satisfy \(q^4<\Lambda\le K_{\mathcal A}\le K_{\mathcal B}\),
\(K_{\mathcal B}/K_{\mathcal A}\le P/\Lambda\le1+1/q\) and
\(K_{\mathcal B}\le kP\le k(q+1)^4\). Put \(x=K_{\mathcal A}^{1/5}-1\)
(\(>q^{4/5}-1>160\)). By Nagura's theorem {[}Na52{]} (for \(y\ge25\)
there is a prime in \((y,\frac65y)\)) there is a prime \(s\) with
\(x/1.2<s<x\). Then \(K_0(s)<K_0(x)<(x+1)^5=K_{\mathcal A}\), and, with
\(K_0(y)\ge y^5\), \(x^5\ge K_{\mathcal A}(1-5K_{\mathcal A}^{-1/5})\)
(Bernoulli) and \(K_{\mathcal A}^{-1/5}<q^{-4/5}\le0.006\),
\[\frac{K_{\mathcal B}}{K_0(s)}<\frac{1.2^5K_{\mathcal B}}{x^5}\le\frac{2.4884\,(1+1/q)}{1-0.03}<3 .\]
So Lemma 4.1 with \(c=3\) gives an
\(\mathrm{HG}(K_{\mathcal B},K_{\mathcal A};3(s+1))\), and
\(3(s+1)<3(x+1)=3K_{\mathcal A}^{1/5}\le3k^{1/5}(q+1)^{4/5}=3(q+1)\bigl(k/(q+1)\bigr)^{1/5}\le 3\cdot300^{-1/5}(q+1) <0.959\,(q+1)\le q-k\),
the last step because \(0.041(q+1)\ge 12.3k\ge k+1\). (b) For
\(q\le200\) the program \texttt{d7thm/hubsum.py} writes for each
\((q,k)\) an explicit decomposition into parts of Lemmas 4.1--4.3,
combined by Lemma 4.2; each line of \texttt{d7thm/hubsum.out} can be
checked by hand (sizes add up, the order condition, degrees). Examples:
\(q=16\), \(k=2\), \(K=66{,}065\): \(H(8)+H(7)+H(5)+H(5)+\) a
pendant-tree part with \(1{,}196\) vertices per colour of degree \(10\);
degrees \(9+4,\ 8+4,\ 6+4,\ 6+4,\ 10+4\le14\). \(q=23\), \(k=3\):
\(H(11)+H(9)+H(8)\) blown up by \(2\) to \(49{,}506\) and \(37{,}867\).
For \(200<q\le2\cdot10^5\) Lemma 4.1 alone suffices
(\texttt{d7thm/hubcap.py}, which finds \(s\) and \(c\) for every such
\(q\)); for \(q>2\cdot10^5\ge300k\) use (a). \(\square\)

\emph{Computer-checked.} \texttt{d7thm/build\_cert.py} builds the hub
graph of a certificate (hexagons from \texttt{research/d10D7/hexgen2.c},
random injections) and \texttt{hgcheck\ edges} checks it:
\((q,k)=(16,2)\), \(\mathrm{HG}(66065,66065;14)\) (the one behind
\((19,7)\)), \((17,2)\), \(\mathrm{HG}(84117,84117;15)\), and
\((23,3)\), \(\mathrm{HG}(293092,281453;20)\) CERT233: no same-colour
pair at distance more than \(6\), degrees as claimed
(\texttt{d7thm/cert\_checks.out}).

\textbf{Computer-checked hub graphs below \(q_k\)}
(\texttt{d7thm/hgcheck.c}, own code: bit-parallel breadth-first search
from every hub; \texttt{d7thm/hgcheck\_lab.out},
\texttt{d7thm/hgrun.out}). The lab's stored hub graphs, read from the
solution files, re-checked standalone: \(\mathrm{HG}(15541,15541;9)\)
{[}(14,7){]}, \(\mathrm{HG}(29535,29535;11)\) {[}(16,7){]},
\(\mathrm{HG}(32356,30327;10)\) and \(\mathrm{HG}(32356,30158;10)\)
{[}(17,7){]}, \(\mathrm{HG}(25308,25308;9)\) {[}(18,7), PSL(2,37){]},
\(\mathrm{HG}(66065,66065;14)\) {[}(19,7){]},
\(\mathrm{HG}(70163,66322;13)\) {[}(20,7){]}: all VALID. The two
cluster-assisted ones {[}(15,7) \(16{,}152\), (18,7) \(34{,}440\){]} are
not hub graphs in this sense (their same-colour pairs are joined through
the levels; in the hub graph alone \(452{,}256\) and \(654{,}360\)
ordered same-colour pairs are at distance more than \(6\), the second
number as in \texttt{conics2/README.md}); for them (HG) is covered by
the all-pairs checks of the whole graphs. New random graphs at the slot
counts of Theorem 2.2 with \(m=q^2+1\): GAPRESULTS

\textbf{Heuristic.} A random bipartite graph of degree \(\delta\) has
same-colour distances at most \(6\) once \(K\) is well below
\(\delta^6/\log\); the failures at \((q,k)=(11,3)\) (\(\delta=8\),
\(K\approx16{,}000\)) and \((13,4)\) (\(\delta=9\),
\(K\approx58{,}000\)) are where this runs out, and there the lab used
Haar graphs over \(\mathrm{PSL}(2,p)\) and paths through the levels.

\hypertarget{the-cells}{%
\section{5. The cells}\label{the-cells}}

\hypertarget{the-fifteen-built-graphs}{%
\subsection{5.1 The fifteen built
graphs}\label{the-fifteen-built-graphs}}

\texttt{d7thm/cells.py} recomputes, for every graph built in the jobs
``conics'' and ``conics-2'', the slot vectors of §2 from \(q\), the type
string, \(|\mathcal P_S|=P(q)\), \(|R_L|\) (\(q^4+q\) without \(S_0\),
\(q^4+q^2+q\) with it) and \(m\), and the order
\(2(2k+1)N(q)+K_{\mathcal A}+K_{\mathcal B}\) from the hub counts.
\textbf{Computer-checked}: all fifteen agree with the builders
(\texttt{cells.out}): every slot vector equals the one printed in the
make logs on branch \texttt{claude/conics-2-diameter-7-ic0oat}, and
every order equals the \(n\) of the solution file. The current graphs:

\begin{longtable}[]{@{}
  >{\raggedright\arraybackslash}p{(\columnwidth - 14\tabcolsep) * \real{0.0800}}
  >{\raggedright\arraybackslash}p{(\columnwidth - 14\tabcolsep) * \real{0.1067}}
  >{\raggedright\arraybackslash}p{(\columnwidth - 14\tabcolsep) * \real{0.0533}}
  >{\raggedright\arraybackslash}p{(\columnwidth - 14\tabcolsep) * \real{0.1333}}
  >{\raggedright\arraybackslash}p{(\columnwidth - 14\tabcolsep) * \real{0.1333}}
  >{\raggedleft\arraybackslash}p{(\columnwidth - 14\tabcolsep) * \real{0.1200}}
  >{\raggedright\arraybackslash}p{(\columnwidth - 14\tabcolsep) * \real{0.1333}}
  >{\raggedright\arraybackslash}p{(\columnwidth - 14\tabcolsep) * \real{0.2400}}@{}}
\toprule\noalign{}
\begin{minipage}[b]{\linewidth}\raggedright
cell
\end{minipage} & \begin{minipage}[b]{\linewidth}\raggedright
\(q\), types
\end{minipage} & \begin{minipage}[b]{\linewidth}\raggedright
\(m\)
\end{minipage} & \begin{minipage}[b]{\linewidth}\raggedright
hubs per side
\end{minipage} & \begin{minipage}[b]{\linewidth}\raggedright
hub graph
\end{minipage} & \begin{minipage}[b]{\linewidth}\raggedleft
\(n\)
\end{minipage} & \begin{minipage}[b]{\linewidth}\raggedright
all-pairs check
\end{minipage} & \begin{minipage}[b]{\linewidth}\raggedright
the bound follows from
\end{minipage} \\
\midrule\noalign{}
\endhead
\bottomrule\noalign{}
\endlastfoot
(14,7) & 11, \(HDHDH\) & 768 & 15,541 + 15,541 & random & 1,802,642 &
passed & Thm 3.1 + computer (set, hub graph) \\
(15,7) & 11, \(HHDHHDD\) & 768 & 16,152 + 16,152 & Haar, clusters &
2,512,488 & passed & all-pairs check \\
(16,7) & 13, \(HDHDH\) & 792 & 29,535 + 29,535 & random & 4,081,410 &
passed & Thm 3.1 + computer (set, hub graph) \\
(17,7) & 13, \(HHDHHDD\) & 792 & 32,356 + 30,327 & random & 5,693,959 &
passed & Thm 3.1 + computer (set, hub graph) \\
(18,7) & 13, \(HHDHDHDHD\) & 792 & 34,440 + 34,440 & Haar, clusters &
7,309,092 & passed & all-pairs check \\
(19,7) & 16, \(HDHDH\) & 257 & 66,065 + 66,065 & random & 11,316,940 &
9/16 ranges & \textbf{Theorem} (\(=F_2(16)\)) \\
(20,7) & 16, \(HHDHHDD\) & 257 & 70,163 + 66,322 & random & 15,795,219 &
9/24 ranges & Thm 3.1 + computer (hub graph) (\(=F_3(16)\)) \\
\end{longtable}

All use \(\Lambda=q^4+q^2+q\) (with \(S_0\)). Hub graphs: ``random''
means random bipartite of degree \(q-k\) (checked standalone, §4);
``Haar, clusters'' means a Haar graph over
\(\mathbb Z_{673}\rtimes\mathbb Z_{24}\) resp. \(\mathrm{PSL}(2,41)\)
whose same-colour pairs are partly joined through the levels
(\texttt{conics2/README.md} §5). The all-pairs states are those of
\texttt{conics2/README.md} at 11:44 UTC.

At \((19,7)\) and \((20,7)\) the built graphs use exactly the sets of
Theorem 3.1 (\(m=q^2+1\), the elliptic section), so their orders are
\(F_2(16)\) and \(F_3(16)\). For \((19,7)\) every ingredient is proved:
the sets (Corollary 1.8, Proposition 1.9), the repair of every
transition (Theorems B and C), the slot counts (Theorem 2.2) and an
\(\mathrm{HG}(66065,66065;14)\) (Theorem 4.4: \(H(8)+H(7)+H(5)+H(5)\)
and a pendant-tree part, joined by Lemma 4.2). So
\[n_b(19,7)\ \ge\ 11{,}316{,}940\] holds by proof, independently of the
all-pairs check of the built graph (which uses a random hub graph and
was not complete when job ``conics-2'' ended). For \((20,7)\) the hub
graph \(\mathrm{HG}(70163,66322;13)\) is not covered by Theorem 4.4; it
is computer-checked (the builder's graph, re-checked standalone by
\texttt{hgcheck}, and a second random graph, seed 1,
\texttt{hgrun.out}); with the proved equal-sided hub graph of
\(65{,}737\) hubs per colour the bound is \(15{,}790{,}208\). This is a
statement about the theorems, not a record claim: the lab records values
after the coordinator's own check.

\hypertarget{the-cells-delta14dots26}{%
\subsection{\texorpdfstring{5.2 The cells
\(\Delta=14,\dots,26\)}{5.2 The cells \textbackslash Delta=14,\textbackslash dots,26}}\label{the-cells-delta14dots26}}

For each \(\Delta\), the best \((q,k)\) with \(q+1+k=\Delta\); ``level
graphs'' is \(\max 2(2k+1)N(q)\), the 2013 table value for
\(\Delta\le16\) and the lab's beyond-table value for
\(\Delta=17,\dots,20\)
(\texttt{research/frontier/beyond/BEYOND\_TABLE.md}, branch
\texttt{claude/new-cells-beyond-etos3t}). \(F_k(q)\) is Corollary 3.2
with \(m=q^2+1\); its hub graph is ``proof'' (Theorem 4.4), ``computer''
(\texttt{hgcheck}) or ``open''. The last column is the best value of
Theorem 3.1 that is proved outright: full slots where Theorem 4.4
applies, otherwise the largest equal-sided hub graph that Lemmas
4.1--4.3 give below the slots (found by \texttt{d7thm/table.py}; the
decomposition is printed in \texttt{table.out}).

TABLE52

The fully proved values exceed the level-graph values at every cell
\(\Delta=14,\dots,26\). For \(k\ge3\) the lab's notes had no general
statement except for \(q=3^h\) (\texttt{D7\_HEXAGON.md}, Corollary 9.2;
§9 there says none covers \((15,7)\)). The values with computer-found
ingredients (the lab's sets \(m=768,792\) and hub graphs) are the built
values of §5.1 and are larger at \(q=11,13\).

\hypertarget{summary-what-is-new}{%
\section{6. Summary: what is new}\label{summary-what-is-new}}

\begin{longtable}[]{@{}
  >{\raggedright\arraybackslash}p{(\columnwidth - 4\tabcolsep) * \real{0.3333}}
  >{\raggedright\arraybackslash}p{(\columnwidth - 4\tabcolsep) * \real{0.3333}}
  >{\raggedright\arraybackslash}p{(\columnwidth - 4\tabcolsep) * \real{0.3333}}@{}}
\toprule\noalign{}
\begin{minipage}[b]{\linewidth}\raggedright
result
\end{minipage} & \begin{minipage}[b]{\linewidth}\raggedright
label
\end{minipage} & \begin{minipage}[b]{\linewidth}\raggedright
status
\end{minipage} \\
\midrule\noalign{}
\endhead
\bottomrule\noalign{}
\endlastfoot
Theorem 1.4: lines meeting \(v\), whiskers and the spread lines missed
by them form a dual segment set, for any hexagon, spread and whisker
set; maximality criterion & new in this form (special cases:
\texttt{CONICS.md} Lemma 3.1, \texttt{conics2} Prop. 1(c),(d)) &
Theorem, coordinate-free \\
Proposition 1.5: in a dual segment set whose lines meet only in spread
points, at most \(q^4+q\) lines are off the spread;
\(\mathcal L_v\cup W\) attains it & new & Theorem \\
Lemma 1.6: an exact whisker set misses at most \(q^3-q^2\) spread lines,
exactly \(q^2\) under the ``all or nothing'' property & new & Theorem \\
Theorem 1.7, Corollary 1.8: \(\lvert R_H\rvert=(q+1)(q^3+1)\),
\(\lvert R_L\rvert=q^4+q^2+q\), \(R_L\) maximal, every odd \(q\) and
even \(q\ge4\) & restatement of \texttt{conics2} Prop. 1 with proof &
Theorem; Computer-checked \(q=2,\dots,16\) (squad c1--c3, builder) \\
Theorem 2.1: for every regular tournament and typing, the hub counts per
side are at most \(\lfloor k/2\rfloor\) friendly sets, with the exact
second-order term & new & Theorem; Computer-checked (all tournaments on
\(\le7\) levels, SAT to 13 levels, squad c6) \\
Theorem 2.2: the alternating layout attains Theorem 2.1 for every \(k\);
closed forms & new & Theorem; Computer-checked \(k\le40\) \\
Theorem 3.1, Corollary 3.2: \(n_b(q+1+k,7)\ge F_k(q)\) given the hub
graph & new (for \(k=2\) it improves \texttt{D7\_HEXAGON.md} Cor. 9.3
from \(\approx2q^3\) to \(\approx2q^4\) above the level graph; nothing
was stated for \(k\ge3\), \(q\neq3^h\)) & Theorem \\
Lemma 4.2: sums of hub graphs; Lemma 4.1: blow-ups; Lemma 4.3:
\texttt{D7\_HEXAGON.md} Lemma 7.1 with the colours sized separately &
4.2 new; 4.1 elementary; 4.3 restatement & Theorem; Computer-checked
(squad c5) \\
Theorem 4.4: the hub graphs at full slots exist for \(q\ge q_k\)
(\(q_2=16\), \(q_3=q_4=23\), \ldots, \(q_k\le300k\)) & new & Theorem
(certificates by program for \(q\le200\)) \\
\(n_b(19,7)\ge11{,}316{,}940\) by proof; fully proved values above the
level graphs at \((14,7)\)--\((26,7)\) (§5.2) & new consequence &
Theorem \\
The formulas reproduce all fifteen built graphs; the lab's hub graphs
re-checked standalone; new random hub graphs for \(q<q_k\) & --- &
Computer-checked \\
Random hub graphs fail at \((q,k)=(11,3),(13,4)\); larger dual segment
sets than \(q^4+q^2+q\) for \(q\ge5\); a coordinate-free description of
the \(q^2\) lines with \(\nu=0\) & --- & open; Heuristic remarks in §1.3
and §4 \\
\end{longtable}

Nothing here raises a value built by the two jobs: the theorems explain
them (at \(q=16\) exactly) and give the same kind of bound for every
\(q\) and every \(k\), with an explicit hub graph wherever \(q\ge q_k\).

\hypertarget{appendix-files}{%
\section{Appendix: files}\label{appendix-files}}

In \texttt{research/synthetic/d7thm/} (Python 3;
\texttt{pip\ install\ python-sat} for \texttt{slots.py}; C with OpenMP
for \texttt{hgcheck.c}):

\begin{itemize}
\tightlist
\item
  \texttt{slots.py} → \texttt{slots.out}: §2 (all typings of the
  rotational tournament for \(k\le8\), all regular tournaments on 5 and
  7 levels, SAT for 9--13 levels).
\item
  \texttt{hubcap\_lib.py}, \texttt{hubcap.py} → \texttt{hubcap.out}:
  Lemma 4.1 coverage for \(q\le2\cdot10^5\) and the inequalities of
  Theorem 4.4(a); \texttt{hubsum.py} → \texttt{hubsum.out}: the
  certificates of Theorem 4.4(b) for \(q\le200\) (one line per
  \((q,k)\)).
\item
  \texttt{hgcheck.c}: standalone check of a hub graph (same-colour
  distances by bit-parallel BFS from every vertex); modes \texttt{sol}
  (read a solution file), \texttt{rand} (random bipartite graph,
  optional aimed 2-switch repair), \texttt{haar} (Haar graphs of
  \(\mathbb Z_K\)). \texttt{hgcheck\_lab.out}: the lab's hub graphs;
  \texttt{hgrun.py} → \texttt{hgrun.out}: random hub graphs at the slot
  counts of Theorem 2.2 for \(q<q_k\).
\item
  \texttt{cells.py} → \texttt{cells.out}: §5.1 and the closed forms of
  Corollary 3.2; \texttt{table.py} → \texttt{table.out}: §5.2.
\end{itemize}

Squad round \texttt{research/squads/d7thm-r1/}: \texttt{BRIEF.md}, cards
\texttt{c1}--\texttt{c6} (code and logs), \texttt{scoreboard.jsonl};
logged in \texttt{research/squads/LOG.md}. The PDF
\texttt{D7\_THEOREMS.pdf} is made from this file by pandoc and tectonic
(\texttt{d7thm/make\_pdf.sh}).

\hypertarget{references}{%
\section{References}\label{references}}

\begin{itemize}
\tightlist
\item
  {[}MS13{]} M. Miller, J. Širáň, Moore graphs and beyond: a survey of
  the degree/diameter problem, Electron. J. Combin., Dynamic Survey DS14
  (2013).
\item
  {[}Na52{]} J. Nagura, On the interval containing at least one prime
  number, Proc. Japan Acad. 28 (1952) 177--181.
\item
  {[}Th80{]} J. A. Thas, Polar spaces, generalized hexagons and perfect
  codes, J. Combin. Theory Ser. A 29 (1980) 87--93.
\item
  {[}VM98{]} H. Van Maldeghem, Generalized Polygons, Birkhäuser, 1998.
\end{itemize}

\end{document}
